\documentclass[sigconf,nonacm]{acmart}

\settopmatter{printacmref=false}
\renewcommand\footnotetextcopyrightpermission[1]{}
\pagestyle{plain}

\usepackage{booktabs}
\emergencystretch=3em
\hbadness=10000
\sloppy

\begin{document}

\title{Monolithic vs.\ Agentic vs.\ Swarm:\\A Controlled Comparison of Three LLM Execution Architectures}
\subtitle{Draft Results Paper}

\author{Eric Medlock}
\affiliation{%
  \institution{University of North Carolina at Charlotte}
  \department{ITCS 6881 Independent Study}
  \city{Charlotte}
  \state{North Carolina}
  \country{USA}
}
\email{emedlock@charlotte.edu}

\author{Dr.\ Jinzhen Wang}
\affiliation{%
  \institution{University of North Carolina at Charlotte}
  \department{Faculty Advisor}
  \city{Charlotte}
  \state{North Carolina}
  \country{USA}
}

\renewcommand{\shortauthors}{Medlock and Wang}

\begin{abstract}
Large language models (LLMs) can be wired together in many ways to solve a
task. The simplest is to ask the model once and take its answer. More elaborate
designs let one model check another's work, or let several copies vote. These
designs cost more---in time, in tokens, and in hardware---so a natural question
is whether the added machinery actually buys more correct answers, and if so,
under what conditions. This paper reports a controlled comparison of three such
designs---\emph{monolithic}, \emph{agentic}, and \emph{swarm}---run under
identical model, prompt, and decoding settings so that the \emph{architecture}
is the only thing that changes. Across a 15-task suite drawn from three standard
benchmarks (GSM8K, HumanEval, and HotpotQA), the agentic design reached 100\%
accuracy and was the best trade-off between accuracy and cost; the monolithic
design was cheapest; and the swarm design cost the most while matching, but not
beating, the monolithic design. We argue that this last result is an artifact of
the tasks being too easy for voting to help, and we describe a harder ``frontier''
task tier, now under construction, designed to reveal where each architecture
genuinely wins.
\end{abstract}

\keywords{large language models, agentic systems, multi-agent swarm,
benchmarking, Pareto analysis, reproducibility}

\maketitle

\section{Introduction}

A large language model (LLM) is a program that, given some text, predicts
plausible text to follow it. In practice this single capability---continue the
text---can be arranged into very different problem-solving \emph{systems}. One
can call the model a single time and use whatever it produces. One can instead
have the model produce an answer and then call it again to critique and repair
that answer, looping until it is satisfied. Or one can run several independent
copies of the model on the same question and take the majority answer, on the
theory that independent mistakes cancel out. These arrangements---call them
\emph{architectures}---differ enormously in how much time, how many generated
tokens, and how much hardware they consume, and it is not obvious that the
expensive ones are worth it.

This paper studies that question directly. We fix everything except the
architecture: the same underlying model, the same decoding settings (temperature,
context window, token budget), the same tasks, and the same hardware. We then
compare three architectures on four axes that matter in practice---accuracy,
latency (wall-clock time), token consumption, and hardware resource use---and ask
not just \emph{which} architecture is best, but \emph{under what conditions} each
one wins. Our guiding hypothesis is that no single architecture dominates
universally: the right choice depends on how hard the task is.

The remainder of the paper explains the three architectures
(Section~\ref{sec:arch}), describes the benchmark tasks and why they were chosen
(Section~\ref{sec:tasks}), summarizes how runs were executed and scored
(Section~\ref{sec:method}), presents results (Section~\ref{sec:results}),
interprets them (Section~\ref{sec:discussion}), and lays out the remaining work
needed to complete the study (Section~\ref{sec:next}).

\section{The Three Architectures}
\label{sec:arch}

All three architectures use the same underlying LLM; they differ only in how many
times the model is called and how those calls are wired together.

\paragraph{Monolithic.} The baseline. The task is handed to the model in a single
prompt, and its single response is taken as the answer. There is no checking, no
retry, and no second opinion. This is the fastest and cheapest arrangement, and
it is the one most people mean by ``using an LLM.'' It serves as the yardstick
against which the more elaborate designs are measured.

\paragraph{Agentic.} A two-role loop. One instance of the model, the
\emph{Executor}, produces a candidate answer. A second instance, the
\emph{Verifier}, is prompted specifically to check that answer for errors and
either accept it or send it back for revision. If the Verifier is not satisfied,
the Executor tries again with the feedback. The loop is capped (at two passes
here) so it cannot run forever. The intuition is that catching and fixing
mistakes should raise accuracy at the cost of extra model calls and time.

\paragraph{Swarm.} A parallel vote. Three independent copies of the model, given
the same task, each produce an answer on their own with no communication between
them. A simple aggregation step then takes the majority answer. There is no
central controller and no ranking of the peers; the design bets that if the
individual copies make \emph{different} mistakes, the majority will still land on
the correct answer. This is the most token-hungry design, since it multiplies the
work by the number of peers.

Two of the three designs (agentic and swarm) are implemented as small graphs of
model calls; the monolithic design is a single call. Crucially, every
architecture is accessed through the same software interface and the same
underlying model, so any measured difference is attributable to the architecture
rather than to incidental implementation choices.

\section{Benchmark Tasks}
\label{sec:tasks}

To exercise different kinds of reasoning, the task suite draws five items from
each of three widely used benchmarks---15 tasks in total. Each benchmark targets a
distinct skill and comes with a definite correct answer, which lets every run be
graded automatically.

\paragraph{GSM8K (5 items) --- grade-school math word problems.} GSM8K
(``Grade School Math 8K'') is a well-known collection of arithmetic word problems
that require several steps of reasoning---not just recalling a fact, but working
through a short chain of calculations to reach a single numeric answer (for
example, totaling costs, splitting quantities, or computing a rate). These items
test multi-step numerical reasoning. They are graded by \emph{exact numeric
match}: the model's final number either equals the reference answer or it does
not.

\paragraph{HumanEval (5 items) --- Python programming.} HumanEval is a standard
code-generation benchmark in which each item gives a function signature and a
docstring describing what the function should do, and the model must write the
body. What makes HumanEval convenient is that every item ships with hidden
\emph{unit tests}: a solution is correct only if the generated code actually runs
and passes those tests. This is the strictest form of grading in the suite---the
answer is executed, not merely compared---and it tests the model's ability to
produce working, logically correct programs.

\paragraph{HotpotQA (5 items) --- multi-hop question answering.} HotpotQA is a
question-answering benchmark whose questions cannot be answered from a single
fact. Each question requires combining two or more pieces of information---``hops''
---to reach the answer (for example, identifying an entity from one clue and then
looking up a property of that entity). These items test reading comprehension and
the chaining of facts. They are graded by \emph{normalized string match}, which
compares the model's answer to the reference after standardizing capitalization,
spacing, and punctuation so that trivial formatting differences do not count as
wrong.

Together the three benchmarks span numeric reasoning, program synthesis, and
multi-hop factual reasoning, giving a small but deliberately varied picture of
where each architecture helps or hurts. The suite is intentionally compact so
that every task can be run many times under every architecture within a
reasonable time and hardware budget.

\section{Method in Brief}
\label{sec:method}

Every one of the 15 tasks was run under every one of the 3 architectures, and each
of those combinations was repeated $N=5$ times to average out run-to-run
variation---$15 \times 3 \times 5 = 225$ runs in total. All runs used a single
open-weights model, \texttt{deepseek-r1-distill-qwen-14b} (a 14-billion-parameter
model, 4-bit quantized), served locally through LM~Studio on an Apple M5 Max with
48\,GB of unified memory, using the machine's Metal GPU backend. Decoding was held
fixed across all runs (temperature 0, an 8{,}192-token context, and a 6{,}144-token
generation budget) so that the model behaved as deterministically as possible.

Each run recorded not only the answer but also its latency, the number of tokens
consumed, and a snapshot of hardware telemetry (memory and GPU use). Answers were
scored two ways. First, an \emph{automatic grader} appropriate to the benchmark
(exact-numeric, unit-test, or normalized-string, as described above) produced an
objective correct/incorrect verdict. Second, as a cross-check, a separate and
deliberately different model---Llama-3.2-3b, chosen from a different model family
than the backend---acted as an ``LLM-as-judge,'' independently rating each answer.
Using a different-family judge guards against a model simply agreeing with its own
style. The judge's verdicts agreed with the automatic grader 93--100\% of the
time, which gives confidence that the automatic scores are trustworthy.

\section{Results}
\label{sec:results}

Table~\ref{tab:results} summarizes the headline numbers, averaged over all tasks
and repetitions. Accuracy is the fraction of runs graded correct; latency is
mean wall-clock seconds per run (with standard deviation); tokens is the mean
number of tokens generated per run.

\begin{table}[h]
  \centering
  \caption{Results across the 15-task baseline suite (225 runs).
  \textbf{Bold} marks the best value in each column.}
  \label{tab:results}
  \begin{tabular}{lccc}
    \toprule
    \textbf{Architecture} & \textbf{Accuracy} & \textbf{Latency (s)} & \textbf{Tokens} \\
    \midrule
    Monolithic  & $93\%$         & $\mathbf{15.2 \pm 17.5}$ & $\mathbf{832}$ \\
    Agentic     & $\mathbf{100\%}$ & $32.5 \pm 32.0$ & $1961$ \\
    Swarm       & $93\%$         & $32.9 \pm 33.1$ & $2649$ \\
    \bottomrule
  \end{tabular}
\end{table}

Three observations stand out. First, the \textbf{agentic} architecture was the
only one to get every task right (100\%), at roughly double the time and tokens
of the monolithic baseline. Second, the \textbf{monolithic} architecture was by
far the cheapest---fastest and fewest tokens---while still getting 93\% correct.
Third, the \textbf{swarm} architecture was the most expensive of all (most tokens,
comparable latency to agentic) yet matched, and did not beat, the monolithic
baseline's 93\% accuracy.

Put in trade-off terms: agentic and monolithic each sit on the ``Pareto
frontier''---the set of options for which you cannot get more accuracy without
paying more, or pay less without losing accuracy. Agentic is the choice when
accuracy matters most; monolithic is the choice when cost matters most. The swarm,
on this suite, is dominated: it costs more than monolithic for the same accuracy.

\section{Discussion}
\label{sec:discussion}

The swarm result deserves scrutiny, because on its face it suggests majority
voting is simply not worth it. We believe the more accurate reading is that
\emph{these particular tasks are too easy for voting to help}. Majority voting
only rescues a task when the individual copies disagree---some right, some
wrong---so that the correct answer can win the vote. On tasks the model almost
always gets right on the first try, all copies already agree, and the vote merely
re-confirms an answer the cheap monolithic design would have produced anyway. The
extra copies add cost without adding information. In other words, the swarm's poor
showing is a property of the \emph{task selection}, not a verdict on the
architecture.

This directly motivates the central hypothesis of the study: architecture
superiority is \emph{task-dependent}. On easy tasks, the cheapest design that
gets the answer wins, and elaboration is wasted. On genuinely hard tasks---where
the model is right only some of the time---checking (agentic) and voting (swarm)
have room to change the outcome. The baseline suite establishes the easy end of
that spectrum; it cannot, by construction, reveal the hard end.

\section{Project Status}
\label{sec:status}

At the time of writing, the study is at the end of its first phase. Concretely:

\begin{itemize}
  \item \textbf{v1 data collection: complete.} All 225 baseline runs
        (15 tasks $\times$ 3 architectures $\times$ 5 repetitions) have been
        executed, hardware-stamped, and independently scored by the
        LLM-as-judge. The dataset is frozen and will not change.
  \item \textbf{v1 analysis: complete.} Accuracy, latency, token, and
        judge-agreement figures are finalized and reported above; the
        accuracy--cost Pareto picture is settled.
  \item \textbf{Software harness: stable.} The shared backend interface, the
        three architectures, the automatic graders, the telemetry capture, the
        judge, and the analysis pipeline are all implemented and in use, backed
        by an offline test suite.
  \item \textbf{Frontier task tier: designed, not yet run.} A candidate set of
        11 harder items exists but is uncalibrated; this is the active work
        item and the critical path to the study's central result.
  \item \textbf{Cloud/GPU environment and framework migration: designed, not
        started.} Provisioning scripts exist; the multi-agent-framework and
        ensemble-swarm extensions are specified but unimplemented.
\end{itemize}

\noindent In short: the easy end of the difficulty spectrum is fully measured,
and the infrastructure needed to measure the hard end is built. The one thing
standing between the current results and the study's main hypothesis is
calibrating and running the frontier tier.

\section{Next Steps}
\label{sec:next}

\paragraph{Frontier task tier (in progress).} The immediate priority is a harder
task tier deliberately chosen to sit in the regime where architecture matters. A
candidate set of 11 harder items (covering math, code, and multi-hop reasoning)
has been assembled. Each candidate will be calibrated by running the monolithic
design on it several times and measuring how often it succeeds; only items where
the single-pass success rate falls in a middle band---hard enough to fail
sometimes, easy enough to succeed sometimes---will be kept, since those are
exactly the tasks where checking and voting can flip the result. Once this tier is
frozen, all three architectures will be run and scored on it, which is expected to
produce the task-by-architecture picture the baseline suite cannot.

\paragraph{Cloud / GPU comparison.} The current results come entirely from Apple
Metal hardware. Provisioning scripts for an NVIDIA GPU cloud instance already
exist but have not been run; executing them would add a second hardware
environment and let us compare the architectures' resource behavior across very
different accelerators.

\paragraph{Planned extensions.} Two larger extensions are designed but not yet
executed. The first re-implements the agentic and swarm designs on established
multi-agent frameworks (CrewAI and OpenAI Swarm respectively) rather than the
current hand-built versions; because the framework itself affects overhead, those
results would be reported as their own baseline rather than compared directly to
the numbers here. The second adds an \emph{ensemble} swarm---voting across not
just multiple copies of one model but multiple prompt variants and multiple
distinct models---to test whether diversity among voters, rather than mere
repetition, is what makes voting pay off.

\section{Conclusion}
\label{sec:conclusion}

Holding everything but the architecture fixed, we found that on a suite of
moderate-difficulty tasks the agentic ``check-and-repair'' design bought perfect
accuracy at a real but modest cost premium, the plain monolithic design was the
efficient choice, and the parallel-voting swarm design paid more for no accuracy
gain---an outcome we attribute to the tasks being too easy for voting to matter.
The work confirms the qualitative shape of our hypothesis and points squarely at
the missing piece: a harder task tier, now under construction, that should reveal
whether and where the more elaborate architectures earn their cost.

\section*{AI Usage Statement}

AI tools were used during this project in a supporting role, with all research
design, decisions, and interpretation remaining the authors' responsibility.
Specifically, LLM-based coding assistants were used to help scaffold and refactor
the experimental harness, write tests, and format results; and an AI assistant
was used to help draft and copy-edit the prose of this report. The empirical
results themselves were produced by the experimental system described in
Section~\ref{sec:method}, not generated by an AI assistant. Notably, an LLM also
appears \emph{inside} the methodology---as the model under test in all three
architectures, and as the independent LLM-as-judge used to cross-check scoring;
those uses are experimental instruments and are documented as such. All
AI-assisted text and code were reviewed and verified by the authors before
inclusion.

\end{document}
